\documentclass[a4paper,12pt]{article}

% --------- XeLaTeX: шрифты ----------
\usepackage{fontspec}
\usepackage{polyglossia}
\setdefaultlanguage{russian}

\usepackage{iftex}
\ifPDFTeX
  \errmessage{Этот файл нужно компилировать XeLaTeX (или LuaLaTeX), не pdfLaTeX.}
\fi

\IfFontExistsTF{Times New Roman}{
  \setmainfont{Times New Roman}
}{
  \setmainfont{TeX Gyre Termes}
}

\IfFontExistsTF{Courier New}{
  \setmonofont{Courier New}[Scale=MatchLowercase]
}{
  \setmonofont{Courier}[Scale=MatchLowercase]
}

% --------- Оформление страницы ----------
\usepackage{geometry}
\geometry{left=3cm,right=3cm,top=2.5cm,bottom=2.5cm}

\usepackage{setspace}
\onehalfspacing

\usepackage{indentfirst}
\setlength{\parindent}{0.8cm}
\setlength{\parskip}{0pt}

\usepackage{microtype}

% --------- Заголовки ----------
\usepackage{titlesec}
\setcounter{secnumdepth}{0}

\titleformat{\section}
  {\bfseries\fontsize{18}{22}\selectfont}
  {}
  {0pt}
  {}

\usepackage{etoolbox}
\pretocmd{\section}{\clearpage}{}{}

% --------- Содержание ----------
\usepackage{tocloft}
\renewcommand{\contentsname}{Содержание}
\renewcommand{\cftsecfont}{\bfseries}
\renewcommand{\cftsecpagefont}{\bfseries}
\setlength{\cftbeforesecskip}{10pt}

% --------- Таблицы ----------
\usepackage{array}
\usepackage{multirow}
\usepackage{booktabs}
\usepackage{caption}
\captionsetup{labelsep=colon}
\usepackage{float}
\usepackage{longtable}

% Инлайн код
\newcommand{\code}[1]{{\ttfamily #1}}

\begin{document}

% ===================== ТИТУЛЬНЫЙ ЛИСТ =====================
\begin{titlepage}
\begin{center}
Московский государственный университет\\
имени М.\,В.\,Ломоносова\\
Факультет вычислительной математики и кибернетики
\end{center}

\vfill

\begin{center}
{\fontsize{18}{22}\selectfont Отчет по заданию №1}\\[6pt]
{\bfseries\fontsize{18}{22}\selectfont «Методы сортировки»}\\[10pt]
{\bfseries\fontsize{16}{20}\selectfont Вариант 7}
\end{center}

\vfill

\begin{flushright}
Выполнила:\\
студентка 104 группы\\
Конгаа А.\,Н.\\[14pt]
Преподаватель:\\
Гуляев Д.\,А.
\end{flushright}

\vfill

\begin{center}
Москва\\
2026
\end{center}
\end{titlepage}

% ===================== СОДЕРЖАНИЕ =====================
\tableofcontents

% ===================== ПОСТАНОВКА ЗАДАЧИ =====================
\section{Постановка задачи}

Требуется реализовать два метода сортировки массива чисел и провести их
экспериментальное сравнение по числу сравнений и перемещений элементов.

Вариант 7 задает следующие параметры:
\begin{itemize}
\item размеры массива: $n \in \{10, 100, 1000, 10000\}$;
\item тип элементов массива: \code{double};
\item порядок сортировки: по возрастанию;
\item сравниваемые методы: Selection Sort и Merge Sort (bottom-up).
\end{itemize}

Для каждого $n$ генерируются массивы четырех типов (тип 1..4), затем для каждого типа
проводится несколько запусков и усреднение результатов.

Целью работы является экспериментальное сравнение эффективности указанных
алгоритмов по количеству выполняемых сравнений и перемещений элементов.

% ===================== РЕЗУЛЬТАТЫ =====================
\section{Результаты экспериментов}

В работе проведено экспериментальное сравнение сортировки выбором и сортировки слиянием (bottom-up). Для массивов размеров $n = 10, 100, 1000, 10000$ измерялось среднее число сравнений и перемещений элементов (по 5 запусков для каждого случая). Результаты приведены в таблицах 1–2.

\begin{table}[H]
\centering
\small
\renewcommand{\arraystretch}{1.15}
\setlength{\tabcolsep}{6pt}

\resizebox{\textwidth}{}{%
\begin{tabular}{|c|c|c|c|c|c|c|}
\hline
\multirow{2}{*}{$n$} & \multirow{2}{*}{Параметр} &
\multicolumn{4}{c|}{Номер сгенерированного массива} & \multirow{2}{*}{Среднее значение} \\
\cline{3-6}
 & & 1 & 2 & 3 & 4 & \\
\hline
\multirow{2}{*}{10}
& Сравнения & 45 & 45 & 45 & 45 & 45 \\
& Перемещения & 0 & 15 & 20 & 22 & 14 \\
\hline
\multirow{2}{*}{100}
& Сравнения & 4950 & 4950 & 4950 & 4950 & 4950 \\
& Перемещения & 0 & 150 & 282 & 283 & 178 \\
\hline
\multirow{2}{*}{1000}
& Сравнения & 499500 & 499500 & 499500 & 499500 & 499500 \\
& Перемещения & 0 & 1500 & 2955 & 2947 & 1850 \\
\hline
\multirow{2}{*}{10000}
& Сравнения & 49995000 & 49995000 & 49995000 & 49995000 & 49995000 \\
& Перемещения & 0 & 15000 & 29683 & 29664 & 18586 \\
\hline
\end{tabular}}
\caption{Таблица 1: Результаты работы Selection Sort}
\end{table}

\begin{longtable}{|c|c|c|c|c|c|c|}
\hline
\multirow{2}{*}{$n$} & \multirow{2}{*}{Параметр} &
\multicolumn{4}{c|}{Номер сгенерированного массива} & \multirow{2}{*}{Среднее значение} \\
\cline{3-6}
 & & 1 & 2 & 3 & 4 & \\
\hline
\endfirsthead

\hline
\multirow{2}{*}{$n$} & \multirow{2}{*}{Параметр} &
\multicolumn{4}{c|}{Номер сгенерированного массива} & \multirow{2}{*}{Среднее значение} \\
\cline{3-6}
 & & 1 & 2 & 3 & 4 & \\
\hline
\endhead
\multirow{2}{*}{10}
& Сравнения & 21 & 15 & 21 & 24 & 20 \\
& Перемещения & 80 & 80 & 80 & 80 & 80 \\
\hline
\multirow{2}{*}{100}
& Сравнения & 372 & 316 & 558 & 556 & 450 \\
& Перемещения & 1400 & 1400 & 1400 & 1400 & 1400 \\
\hline
\multirow{2}{*}{1000}
& Сравнения & 5052 & 4932 & 8710 & 8693 & 6846 \\
& Перемещения & 20000 & 20000 & 20000 & 20000 & 20000 \\
\hline
\multirow{2}{*}{10000}
& Сравнения & 71712 & 64608 & 123486 & 123456 & 95815 \\
& Перемещения & 280000 & 280000 & 280000 & 280000 & 280000 \\
\hline
\caption{Таблица 2: Результаты работы Merge Sort}
\end{longtable}

Теоретически сортировка выбором выполняет $\frac{n(n-1)}{2}$ сравнений независимо от
входных данных, а сортировка слиянием имеет трудоёмкость порядка $O(n\log n)$~[1].
Экспериментальные данные это подтверждают: для сортировки выбором число сравнений
точно совпадает с формулой (45, 4950, 499500, 49995000 для указанных $n$), что
отражает квадратичный рост трудоёмкости.

Для сортировки слиянием число сравнений растёт значительно медленнее (примерно
20, 450, 6846, 95815), что соответствует асимптотике $O(n\log n)$. При увеличении
размера массива в 10 раз число операций возрастает примерно в 10–15 раз, тогда
как у сортировки выбором — примерно в 100 раз.

Таким образом, при малых $n$ различия между алгоритмами невелики, однако при
больших размерах массивов сортировка слиянием существенно эффективнее сортировки
выбором по числу сравнений и перемещений.

% ===================== СТРУКТУРА =====================
\section{Структура программы и спецификация функций}

Программа реализована на языке C и имеет модульную структуру. Каждый модуль
отвечает за отдельную функциональную часть: генерацию исходных данных,
реализацию алгоритмов сортировки и проведение экспериментальных исследований.

В состав программы входят следующие файлы: \code{main.c}, \code{fill.c},
\code{selection.c}, \code{merge\_bottom\_up.c} и заголовочный файл \code{sort.h}.\\

Заголовочный файл \code{sort.h} содержит объявления глобальных счётчиков операций и
прототипы всех функций программы:
\begin{itemize}
\item \code{extern long long cmp\_count;} — глобальный счётчик числа сравнений;
\item \code{extern long long mov\_count;} — глобальный счётчик числа перемещений;
\item \code{void selection\_sort(double* a, int n);} — функция сортировки выбором;
\item \code{void merge\_sort\_bottom\_up(double* a, int n);} — функция итеративной сортировки слиянием (bottom-up);
\item \code{void fill\_array(double* a, int n, int type);} — функция генерации тестовых массивов;
\item \code{void copy\_array(double* dst, double* src, int n);} — функция копирования массива.
\end{itemize}\\

Файл \code{fill.c} содержит функции формирования исходных данных.
Функция \code{fill\_array} заполняет массив одним из четырёх способов:
\begin{enumerate}
\item упорядоченная последовательность: $a[i] = i$;
\item обратный порядок: $a[i] = n - i$;
\item случайные значения в диапазоне $[0, 99]$;
\item независимая генерация случайного массива.
\end{enumerate}

Функция \code{copy\_array} выполняет копирование массива и используется для
сохранения исходных данных перед выполнением сортировки.\\

Файл \code{selection.c} содержит реализацию алгоритма сортировки выбором —
функцию \code{selection\_sort}. Алгоритм выполняет упорядочивание массива по
возрастанию и в процессе работы подсчитывает количество операций:
\begin{itemize}
\item сравнений (\code{cmp\_count});
\item перемещений (\code{mov\_count}).
\end{itemize}\\

Файл \code{merge\_bottom\_up.c} содержит реализацию итеративной сортировки
слиянием — функцию \code{merge\_sort\_bottom\_up}. Алгоритм выполняет
последовательное слияние соседних отсортированных блоков массива с
постепенно увеличивающимся размером.

Для объединения двух отсортированных участков используется вспомогательная
функция \code{merge}, которая применяет временный массив для корректного
слияния элементов с сохранением их порядка. В процессе выполнения также
ведётся подсчёт числа сравнений и перемещений элементов.\\

Файл \code{main.c} содержит функцию \code{main}, организующую проведение
вычислительных экспериментов.

В функции:
\begin{itemize}
\item задаются размеры массивов и типы исходных данных;
\item выполняется генерация тестовых массивов;
\item вызываются алгоритмы сортировки;
\item производится многократный запуск с последующим усреднением значений счётчиков операций;
\item результаты выводятся в табличной форме.
\end{itemize}

% ===================== ТЕСТЫ =====================
\section{Отладка программы, тестирование функций}

В процессе разработки выполнялись отладка и тестирование реализованных
алгоритмов сортировки. Проверялась корректность сортировки (выходной
массив должен быть упорядочен по возрастанию) и корректность подсчёта
числа сравнений и перемещений.

Тестирование проводилось на массивах разных типов: случайных, уже
отсортированных и упорядоченных в обратном порядке. Использовались
контрольные тесты:
\begin{itemize}
\item Тест 1: $n = 5$, \{3, 1, 2, 2, 0\} $\rightarrow$ \{0, 1, 2, 2, 3\}.
\item Тест 2: тип 1, $n = 10$ (отсортированный массив) $\rightarrow$ без изменений.
\item Тест 3: тип 2, $n = 10$ (обратный порядок) $\rightarrow$ сортировка по возрастанию.
\end{itemize}

Отладка выполнялась с помощью пошагового выполнения и анализа
промежуточных состояний массивов. Особое внимание уделялось границам
подмассивов в сортировке слиянием и выбору минимального элемента в
сортировке выбором.

Тестирование подтвердило корректную работу функций сортировки для всех
исследованных типов данных и размеров массивов.

% ===================== ОШИБКИ =====================
\section{Анализ допущенных ошибок}

Основные ошибки были связаны с реализацией сортировки слиянием
(bottom-up): некорректная обработка границ подмассивов могла приводить
к выходу за пределы массива. После уточнения условий циклов и проверки
индексов ошибка устранена.

Также была уточнена методика подсчёта операций (перемещение — запись в
массив или буфер) и добавлено освобождение временного массива.
Реализация сортировки выбором ошибок не содержала. Тестирование
подтвердило корректность обеих реализаций.

% ===================== ЛИТЕРАТУРА =====================
\section{Список литературы}

[1] Седжвик Р. Алгоритмы на Java. В 2 ч. Ч. 1–2. 4-е изд. — М.: Вильямс, 2011.

\end{document}